% =====================================================================
% 1b. İLGİLİ ÇALIŞMALAR (Related Work)
% Qualitative survey of prior work; describes WHAT others did. The numeric
% comparison to our results stays in §4.9 (subsec:contextual) and is only
% cross-referenced here. Cite ONLY paper.md-backed bib entries. No numbers,
% no SOTA claims.
% =====================================================================
\section{İlgili Çalışmalar}
\label{sec:related}

Gastrointestinal endoskopi görüntülerinin sınıflandırılması, HiperKvasir veri
kümesinin~\cite{borgli2020hyperkvasir} yayımlanmasından bu yana çeşitli
çalışmalarda ele alınmıştır. Borgli ve ark.~\cite{borgli2020hyperkvasir} veri
kümesini tanıtmış ve resmi bölmeler üzerinde önceden eğitilmiş evrişimsel ağlarla
(ResNet, DenseNet ve bunların topluluğu) baseline sonuçları raporlamıştır.
Ramamurthy ve ark.~\cite{ramamurthy2022effimix} (Effimix), EfficientNet tabanlı
çok-özellikli bir füzyon yaklaşımı önermiş; ancak değerlendirmeyi veri artırma ile
dengelenmiş ve resmi olmayan bir bölme üzerinde yapmıştır. GastroViT
\cite{gastrovit2025} görüntü dönüştürücü (Vision Transformer) tabanlı bir topluluk
kullanmış; Ahmed ve ark.~\cite{ahmed2023hybrid} ise derin özellikler ile klasik
sınıflandırıcıları birleştiren hibrit bir füzyon şeması sunmuştur. Bu çalışmalar
farklı veri bölme ve ön işleme protokolleri kullandığından sonuçlarımızla doğrudan
baş-başa karşılaştırılamaz; bu nedenle sayısal kıyaslama yalnızca
Bölüm~\ref{subsec:contextual}'da, protokol farkları açıkça belirtilerek bağlamsal
biçimde ele alınmıştır.

Özellik füzyonu için literatürde farklı stratejiler bulunur. En yalın yaklaşım,
dal özelliklerini birleştirme (concatenation) ile uç uca eklemektir; öğrenilebilir
ağırlıklı füzyon ise dalların katkısını eğitim sırasında ayarlanabilir kılar. Daha
gelişmiş bir alternatif, her örnek için kapı (gate) ağırlıkları hesaplayan kapılı
çok-kipli birimdir (Gated Multimodal Unit, GMU)~\cite{arevalo2017gmu}. Bu
çalışmada, zorunlu iki yöntem olan birleştirme ve ağırlıklı füzyonun yanı sıra GMU
varyantı da uygulanmış ve aynı protokol altında karşılaştırılmıştır.

Sınıf dengesizliği, dengesiz görsel sınıflandırma görevlerinde yaygın bir zorluktur
ve çoğunlukla kayıp düzeyinde ele alınır. Lin ve ark.'nın focal
loss~\cite{lin2017focal} yaklaşımı, kolay örneklerin katkısını azaltarak eğitimi
zor örneklere yöneltir. Bu çalışmada focal loss, en iyi yapılandırma üzerinde
tek-değişkenli bir ablasyon olarak denenmiş ve sonucu Bölüm~\ref{sec:results}'te
verilmiştir.

Bu çalışma, anılan yaklaşımlardan farklı olarak, üç önceden eğitilmiş omurganın
(ResNet50, MobileNetV2, EfficientNetB0) ağırlıklı füzyonu ile MLP sınıflandırmasını
resmi 5-katlı protokol altında ve önyükleme güven aralıklarıyla değerlendiren
kontrollü bir karşılaştırma sunar.
